 % Template:     Auxiliar LaTeX
% Documento:    Archivo principal
% Versión:      8.3.4 (07/02/2025)
% Codificación: UTF-8
%
% Autor: Pablo Pizarro R.
%        pablo@ppizarror.com
%
% Manual template: [https://latex.ppizarror.com/auxiliares]
% Licencia MIT:    [https://opensource.org/licenses/MIT]

% CREACIÓN DEL DOCUMENTO
\documentclass[
	spanish, % Idioma: spanish, english, etc.
	letterpaper, oneside
]{article}


\usepackage{amsmath}
\usepackage{breqn}

\newif\ifshowcomments
\showcommentstrue % Set to \showcommentsfalse to hide all comments

\usepackage{xcolor} % For colored text
\usepackage{marginnote} % Optional: for margin notes

% Define inline comment
\newcommand{\comment}[1]{%
  \ifshowcomments
    {\textcolor{red}{\textbf{ Solución:} }\\ #1 }%
  \fi
}
 

% INFORMACIÓN DEL DOCUMENTO
\def\documenttitle {Actividad Complementaria RP}
\def\documentsubject{Optimización bajo incertidumbre}

\def\documentauthor {Joaquín Ormazábal}
\def\coursename {Optimización}
\def\coursecode {EL4114}

\def\universityname {Universidad de Chile}
\def\universityfaculty {Facultad de Ciencias Físicas y Matemáticas}
\def\universitydepartment {Departamento de Ingeniería Eléctrica}
\def\universitydepartmentimage {departamentos/die}
\def\universitydepartmentimagecfg {height=1.75cm}
\def\universitylocation {Santiago de Chile}

% EQUIPO DOCENTE
\def\teachingstaff {
	\textbf{Profesor: Heraldo Rozas} \\
	Auxiliares: Diego Carez, Joaquín Ormazábal
}

% IMPORTACIÓN DEL TEMPLATE
\input{template}

% INICIO DE PÁGINAS
\begin{document}

% CONFIGURACIÓN DE PÁGINA Y ENCABEZADOS
\templatePagecfg


\fbox{\begin{minipage}{\textwidth}
\noindent\textbf{Instrucciones:} La actividad se realiza de forma individual, con la posibilidad de colaborar en grupos. Trabaje sobre el Jupyter Notebook adjunto en la descripción de la actividad.
\end{minipage}}



% ======================= INICIO DEL DOCUMENTO =======================
\section{Contexto y Motivación}

La creciente integración de fuentes de energía no despachables, como las plantas eólicas y térmico-solares, ha situado a la incertidumbre en el centro de la planificación y operación de los sistemas eléctricos modernos. La generación eólica, por ejemplo, es intrínsecamente variable e incierta, lo que exige que el sistema opere con un grado mucho mayor de flexibilidad para poder adaptarse a sus fluctuaciones. Además de la intermitencia climática, los sistemas eléctricos contemporáneos deben hacer frente a múltiples incertidumbres simultáneas, tales como la volatilidad en la demanda de los clientes, las fallas inesperadas en las unidades de generación y las salidas de servicio o contingencias en las líneas de transmisión\footnote{A. J. Conejo, M. Carrión, and J. M. Morales, \textit{Decision making under uncertainty in electricity markets}, 2010. doi: 10.1007/978-1-4419-7421-1.}.

Abordar estos problemas ignorando la incertidumbre —por ejemplo, sustituyendo las variables aleatorias por sus valores esperados en modelos deterministas— resulta en modelos frágiles y soluciones subóptimas que, al implementarse en el mundo real, fallan en alcanzar los mejores resultados. Asimismo, utilizar enfoques de seguridad puramente deterministas (como prepararse solo para la falla más grande) hace que el sistema subestime los riesgos si no considera las probabilidades reales$^{1,}$\footnote{L. Wu, M. Shahidehpour, and T. Li, ``Stochastic Security-Constrained unit commitment,'' \textit{IEEE Transactions on Power Systems}, vol. 22, no. 2, pp. 800–811, May 2007, doi: 10.1109/tpwrs.2007.894843.}.

Un ejemplo concreto de mala planificación en sistemas de potencia es la crisis de California en agosto de 2020\footnote{N. Groom, ``California blames blackouts on poor planning for extreme heat,'' Reuters, Oct. 2020.}, donde una ola de calor extrema, junto con una menor disponibilidad renovable en horas críticas, provocó cortes rotativos de suministro. Otro caso relevante es Texas en febrero de 2021, cuando la tormenta invernal Uri generó fallas simultáneas en unidades de generación y problemas de abastecimiento de gas natural\footnote{B. Brooks, ``Ice everywhere as Texans brave cold, 2.7 million homes lack power,'' Reuters, Feb. 2021.}. En Chile, la crisis eléctrica de 1998–1999 también evidencia este problema, ya que una sequía severa redujo la disponibilidad hidroeléctrica y llevó a restricciones de suministro\footnote{A. Astudillo M., ``Crisis del gas y racionamiento, entre los puntos bajos de la red eléctrica en los últimos 15 años,'' Diario Financiero, Oct. 2014.}.

Por ello, la \textbf{optimización estocástica} proporciona un marco matemático fundamental\footnote{N. V. Sahinidis, ``Optimization under uncertainty: state-of-the-art and opportunities,'' \textit{Computers \& Chemical Engineering}, vol. 28, no. 6–7, pp. 971–983, Nov. 2003.} para derivar decisiones informadas y robustas frente a la constante incertidumbre del sector.

\subsection{Unit Commitment bajo incertidumbre}

El problema de \textit{Unit Commitment} (UC) consiste en determinar qué unidades generadoras deben estar encendidas o apagadas en cada período de operación, junto con su nivel de generación, con el objetivo de satisfacer la demanda al mínimo costo operativo posible. Este problema es naturalmente MILP debido a la presencia de variables binarias que representan el estado de encendido de cada unidad.

En una formulación determinista clásica, se asume que la demanda y la disponibilidad de generación renovable son conocidas con certeza. Sin embargo, en sistemas eléctricos con alta penetración de energías renovables, esta hipótesis puede ser restrictiva. 

Históricamente, la planificación y operación de los sistemas eléctricos han dependido en gran medida de modelos de optimización deterministas; en el problema de UC, la incertidumbre se maneja habitualmente resolviendo modelos deterministas con requisitos de reserva conservadores. Aunque este paradigma es computacionalmente sencillo, es económicamente ineficiente porque no modela explícitamente la naturaleza estocástica de la incertidumbre, lo que puede resultar en programaciones operativas subóptimas o estados infactibles cuando las condiciones en tiempo real se desvían de los pronósticos nominales.

En términos matemáticos, plantear el modelo determinista con requisitos de reserva implica incorporar restricciones adicionales de desigualdad que aseguren que la capacidad máxima combinada de las unidades en línea en cada período sea mayor o menor que la demanda pronosticada por un margen de seguridad predefinido (ej. un porcentaje de la carga máxima). De este modo, las variables binarias de acoplamiento temporal y los límites de generación nominal se ven forzados a mantener un ``colchón extra'' de potencia activa disponible.

\subsection{Programación estocástica}

%En numerosos problemas de planificación y operación de sistemas eléctricos, las decisiones deben tomarse antes de conocer con certeza la realización de ciertas variables exógenas. En particular, la demanda eléctrica y la generación renovable presentan variabilidad temporal y errores de pronóstico que pueden afectar significativamente la operación económica y segura del sistema.

La \textit{programación estocástica} proporciona un marco matemático para incorporar explícitamente las fuentes de incertidumbre discutidas en el problema de UC,  mediante un conjunto finito de escenarios. Cada escenario representa una posible realización futura de las variables inciertas y se asocia a una probabilidad de ocurrencia. En esta actividad, usted trabajará con incertidumbre asociada a la demanda eléctrica y a la disponibilidad de generación solar.

En términos generales, un problema de programación estocástica de dos etapas puede escribirse como:

\begin{equation}
\min_{x \in \mathcal{X}}
c^\top x + \mathbb{E}_{\xi} \left[ Q(x,\xi) \right],
\label{eq:fo}
\end{equation}

donde $x$ representa las decisiones de primera etapa, tomadas antes de observar la incertidumbre $\xi$, mientras que $Q(x,\xi)$ corresponde al valor óptimo del problema de segunda etapa una vez revelada la realización de $\xi$.

Cuando la distribución de $\xi$ se aproxima mediante un conjunto finito de escenarios $\mathcal{S}$, con probabilidades $\pi_s$ para cada $s \in \mathcal{S}$, la formulación anterior puede aproximarse mediante el equivalente determinista:
\begin{equation}
\min_{x \in \mathcal{X}}
c^\top x + \sum_{s \in \mathcal{S}} \pi_s Q(x,\xi_s)
\label{eq}
\end{equation}

En el contexto de esta actividad, las decisiones de primera etapa estarán asociadas al compromiso de unidades generadoras (variables binarias de encendido, apagado y status), mientras que las decisiones de segunda etapa estarán asociadas al despacho económico bajo cada escenario de demanda y generación solar. 

Esta actividad busca conectar conceptos teóricos de optimización bajo incertidumbre con una aplicación concreta en sistemas eléctricos. A través de ella, usted podrá observar la importancia de distinguir entre decisiones anticipativas y decisiones adaptativas, así como comprender el efecto que tiene la incertidumbre sobre la planificación operacional de sistemas con generación renovable.

%\begin{comment}
\section{Problema de \textit{Unit Commitment}}

El problema de pre-despacho económico o \textit{unit commitment} consiste en determinar qué unidades generadoras deben estar encendidas o apagadas en cada período de operación, junto con su nivel de generación, con el objetivo de satisfacer la demanda al mínimo costo operativo posible. Este problema es naturalmente MILP debido a la presencia de variables binarias que representan el estado de encendido de cada unidad.

En una formulación determinista clásica, se asume que la demanda y la disponibilidad de generación renovable son conocidas con certeza. Sin embargo, en sistemas eléctricos con alta penetración de energías renovables, esta hipótesis puede ser restrictiva. Por esta razón, en esta actividad usted trabajará con una extensión estocástica del problema, donde la demanda y la generación solar dependen de un escenario $s \in \mathcal{S}$.\\



\begin{comment}
%\section{Unit Commitment bajo incertidumbre}

%El problema de \textit{unit commitment} consiste en determinar qué unidades generadoras deben estar encendidas o apagadas en cada período de operación, junto con su nivel de generación, con el objetivo de satisfacer la demanda al mínimo costo operativo posible. Este problema es naturalmente MILP debido a la presencia de variables binarias que representan el estado de encendido de cada unidad.

%En una formulación determinista clásica, se asume que la demanda y la disponibilidad de generación renovable son conocidas con certeza. Sin embargo, en sistemas eléctricos con alta penetración de energías renovables, esta hipótesis puede ser restrictiva. Por esta razón, 

%Tal vez colocar otra transición
En esta actividad usted trabajará con una extensión estocástica del problema de UC, donde la demanda y la generación solar dependen de un escenario $s \in \mathcal{S}$.\\

Sean los siguientes conjuntos y parámetros del problema:

\begin{itemize}
\item $\mathcal{G}$: conjunto de unidades generadoras convencionales.
\item $\mathcal{T}$: conjunto de períodos de operación.
\item $\mathcal{S}$: conjunto de escenarios.
\item $\pi_s$: probabilidad asociada al escenario $s \in \mathcal{S}$.
\item $D_{t,s}$: demanda eléctrica en el período $t$ bajo el escenario $s$.
\item $\overline{P}^{\text{sol}}_{t,s}$: generación solar máxima disponible en el período $t$ bajo el escenario $s$.
\end{itemize}

Las variables de decisión del modelo son:
\begin{itemize}
\item $u_{g,t} \in {0,1}$: variable binaria igual a $1$ si la unidad $g$ está encendida en el período $t$.
\item $v_{g,t} \in {0,1}$: variable binaria igual a $1$ si la unidad $g$ se enciende al inicio del período $t$.
\item $w_{g,t} \in {0,1}$: variable binaria igual a $1$ si la unidad $g$ se apaga al inicio del período $t$.
\item $p_{g,t,s} \geq 0$: potencia generada por la unidad convencional $g$ en el período $t$ bajo el escenario $s$.
\item $p^{\text{sol}}{t,s} \geq 0$: generación solar utilizada en el período $t$ bajo el escenario $s$.
\item $\ell{t,s} \geq 0$: energía no suministrada en el período $t$ bajo el escenario $s$.
\end{itemize}

La función objetivo minimiza la suma entre costos deterministas de compromiso de unidades y costos esperados de operación. Las restricciones aseguran el balance de potencia para cada período y escenario, limitan la generación convencional y solar según su disponibilidad, y representan la lógica temporal de encendido y apagado de las unidades.
\end{comment}

\begin{comment}
\section{Modelamiento e implementación}
En el mundo laboral utilizarán herramientas de IA para ser más eficientes. Por ello, en esta entrega se fomenta su uso estratégico. El objetivo es que deleguen en la IA las tareas repetitivas de formato o sintaxis, mientras ustedes se concentran en el verdadero núcleo del problema: el modelamiento y el análisis de los resultados.

Es importante que antes de escribir cualquier línea de código, deben estructurar matemáticamente el problema en papel o en su editor de texto científico. La regla principal es si decides usar IA para generar un fragmento de código, debes ser capaz de explicar línea por línea qué hace esa función en una eventual defensa o corrección.

\subsection{Guía para el uso}
Para el desarrollo de esta actividad, se permite y fomenta el uso de herramientas de IA (como ChatGPT, Claude, Gemini, etc.) como \textbf{asistentes de aprendizaje y co-programadores}, bajo las siguientes pautas:

\subsubsection*{Usos Recomendados}
\begin{itemize}[label=\checkmark]
    \item \textbf{Tutor de Sintaxis:} Si tu formulación en papel es correcta, pero olvidas cómo declarar variables indexadas por múltiples conjuntos en tu librería de optimización, pídele a la IA ejemplos genéricos de sintaxis.
    \item \textbf{Explicación de Errores:} Si el optimizador arroja un error de código o de factibilidad (\textit{Infeasible}), copia el error en la IA y pídele que te ayude a interpretar qué restricción o lógica matemática podría estar fallando.
    \item \textbf{Optimización y Limpieza de Código:} Utiliza la IA para mejorar la legibilidad de tus gráficos o para estructurar de manera más limpia tus bucles e impresiones en consola.
\end{itemize}

\subsubsection*{Prácticas no permitidas}
\begin{itemize}[label=$\times$]
    \item \textbf{Copiado a Ciegas (\textit{Black Box}):} Evita pedirle a la IA que escriba el modelo completo desde cero. El código generado por IA suele omitir restricciones sutiles o inventar sintaxis que no corresponde a la librería, lo que te costará más tiempo depurar que escribirlo tú mismo.
    \item \textbf{Delegar el Pensamiento Crítico:} La IA no debe responder las preguntas de discusión por ti. Las herramientas de IA suelen redactar análisis genéricos y superficiales. Tu evaluación debe fundamentarse estrictamente en los datos, números y gráficos que arrojó \textbf{tu} simulación.
\end{itemize}
\end{comment}

\section{Descripción de la actividad}

El objetivo de esta actividad es que usted se familiarice con los conceptos fundamentales de programación estocástica mediante el modelamiento e implementación computacional de un problema simplificado de unit commitment con incertidumbre en demanda y generación solar.

Para ello, deberá trabajar sobre el Jupyter Notebook disponible en el siguiente \href{https://github.com/paoptlab/Apunte-Optimizacion-A2IC}{link}, identificando la estructura matemática del problema, implementando las variables y restricciones correspondientes, y analizando cómo cambian las decisiones óptimas al incorporar escenarios de incertidumbre.

%La actividad se divide en las siguientes etapas:

\subsection{Modelamiento}

Antes de realizar cualquier implementación computacional, debe estructurar analíticamente el problema de UC. El objetivo de esta parte es que comprenda cómo la incorporación de requisitos de reserva del sistema altera las decisiones de despacho y compromiso de las unidades.

\subsection*{A) Modelo determinista}

En primer lugar, considere la versión determinista del problema, en la cual se considera un único pronóstico de demanda y generación solar. Se estudia un problema uninodal, con un planta de generación solar, una punto de demanda y $g$ generadores convencionales, los conjuntos, parámetros y variables del modelo están en la Tabla \ref{tab:notacion_det}. A continuación el desglose de la formulación determinista del problema de UC.


\begin{table}[h]
\centering
\renewcommand{\arraystretch}{1.3}
\begin{tabular}{c l}
\toprule
\textbf{Símbolo} & \textbf{Descripción} \\
\midrule
\multicolumn{2}{l}{\textbf{\textit{Conjuntos e índices}}} \\
\midrule
$\mathcal{G}$ & Conjunto de unidades generadoras convencionales (índice $g$) \\
$\mathcal{T}$ & Conjunto de períodos de operación (índice $t$) \\
\midrule
\multicolumn{2}{l}{\textbf{\textit{Parámetros}}} \\
\midrule
$D_{t}$ & Demanda eléctrica en el período $t$ \\
$\overline{P}^{\text{sol}}_{t}$ & Generación solar máxima disponible en el período $t$ \\
$c_{g}^{\text{var}}$ & Costo variable de generación de la unidad $g$ \\
$c_{g}^{\text{start}}$ & Costo de encendido (arranque) de la unidad $g$ \\
$\text{VOLL}$ & Costo de energía no suministrada (desprendimiento de carga) \\
\midrule
\multicolumn{2}{l}{\textbf{\textit{Variables}}} \\
\midrule
$u_{g,t}\in\{0,1\}$ & $=1$ si la unidad $g$ está encendida en el período $t$ \\
$v_{g,t}\in\{0,1\}$ & $=1$ si la unidad $g$ se enciende al inicio del período $t$ \\
$w_{g,t}\in\{0,1\}$ & $=1$ si la unidad $g$ se apaga al inicio del período $t$ \\
$p_{g,t}\geq 0$ & Potencia generada por la unidad $g$ en el período $t$ \\
$p^{\text{sol}}_{t}\geq 0$ & Generación solar utilizada en el período $t$ \\
$\ell_{t}\geq 0$ & Energía no suministrada en el período $t$ \\
\bottomrule
\end{tabular}
\caption{Conjuntos, parámetros y variables del modelo determinista de \textit{unit commitment}.}
\label{tab:notacion_det}
\end{table}


\subsubsection*{Función objetivo}

La función objetivo minimiza el costo total de compromiso y operación:
\begin{equation}
\begin{aligned}
    \min_{u,v,w,p,p^{sol},\ell} \quad
    &\underbrace{
    \sum_{g\in\mathcal{G}}\sum_{t\in\mathcal{T}}
    \left(
        c^{nl}_g u_{g,t}
        +
        c^{su}_g v_{g,t}
    \right)
    }_{\text{costo de compromiso}}
    +
    \underbrace{
    \sum_{t\in\mathcal{T}}
    \left(
        \sum_{g\in\mathcal{G}} c_g p_{g,t}
        +
        \mathrm{VOLL}\,\ell_t
    \right)
    }_{\text{costo de operación}}.
\end{aligned}
\label{eq:objetivo_determinista}
\end{equation}

Sujeto a las siguientes restricciones:

\paragraph{1. Balance de potencia}

\begin{equation}
    \sum_{g\in\mathcal{G}} p_{g,t}
    +
    p^{sol}_t
    +
    \ell_t
    =
    D_t,
    \qquad 
    \forall t\in\mathcal{T}.
    \label{eq:balance_potencia_determinista}
\end{equation}

\paragraph{2. Mínimo técnico}

\begin{equation}
    p_{g,t}
    \geq
    P^{\min}_g u_{g,t},
    \qquad 
    \forall g\in\mathcal{G},\; t\in\mathcal{T}.
    \label{eq:minimo_tecnico_determinista}
\end{equation}

\paragraph{3. Límite superior de generación}

\begin{equation}
    p_{g,t}
    \leq
    P^{\max}_g u_{g,t},
    \qquad 
    \forall g\in\mathcal{G},\; t\in\mathcal{T}.
    \label{eq:limite_superior_generacion_determinista}
\end{equation}

\paragraph{4. Disponibilidad solar}

\begin{equation}
    p^{sol}_t
    \leq
    P^{sol}_t,
    \qquad 
    \forall t\in\mathcal{T}.
    \label{eq:disponibilidad_solar_determinista}
\end{equation}

\paragraph{5. Lógica de transición}

\begin{equation}
    u_{g,t}-u_{g,t-1}
    =
    v_{g,t}-w_{g,t},
    \qquad 
    \forall g\in\mathcal{G},\; t\in\mathcal{T}.
    \label{eq:logica_transicion_determinista}
\end{equation}

\paragraph{6. Exclusividad arranque/parada}

\begin{equation}
    v_{g,t}+w_{g,t}
    \leq
    1,
    \qquad 
    \forall g\in\mathcal{G},\; t\in\mathcal{T}.
    \label{eq:exclusividad_arranque_parada_determinista}
\end{equation}

En la formulación base el balance se satisface de forma exacta para un único pronóstico, sin margen ante desviaciones. Para dotar al sistema de un colchón frente 
a contingencias, tales como la salida de una unidad o errores de pronóstico, se exige una \textbf{reserva} equivalente a una fracción $\alpha$ de la demanda. Esta formulación no introduce nuevas variables. Se modifica únicamente la ecuación 
de balance de potencia, sobredimensionando la demanda por un factor $(1+\alpha)$:

\begin{equation}
    \sum_{g\in\mathcal{G}} p_{g,t}
    +
    p^{sol}_t
    +
    \ell_t
    =
    (1+\alpha)D_t,
    \qquad 
    \forall t\in\mathcal{T}.
    \label{eq:balance_potencia_reserva_implicita}
\end{equation}


\begin{comment}
\subsection*{Conjuntos}

\begin{itemize}
    \item $T = \{1,\ldots,12\}$: conjunto de índices asociados a cada hora en análisis en el problema.
    \item $G = \{1,\ldots,4\}$: conjunto de índices asociados a cada generador convencional del sistema.
    \item $\mathcal{R} = \{1,\ldots,1\}$: conjunto de índices asociados a cada generador renovable del sistema.
    \item $N = \{1,\ldots,4\}$: conjunto de índices de cada nodo del sistema.
    \item $L = \{1,\ldots,4\}$: conjunto de índices de cada línea de transmisión del sistema.
    \item $G_n$: conjunto de índices de generadores convencionales instalados en el nodo $n$.
    \item $R_n$: conjunto de índices de generadores renovables instalados en el nodo $n$.
\end{itemize}

\subsection*{Parámetros}

\begin{itemize}
    \item $to(l)$: indica el nodo $n$ hacia el cual la línea $l$ entra.
    \item $from(l)$: indica el nodo $n$ desde el cual la línea $l$ sale.
    \item $C_g$: costo variable del generador $g \in G$.
    \item $C^E_g$: costo de encendido del generador $g \in G$.
    \item $C^A_g$: costo de apagado del generador $g \in G$.
    \item $\overline{P}_g$: potencia máxima del generador $g \in G$ cuando está encendido.
    \item $\underline{P}_g$: potencia mínima del generador $g \in G$ cuando está encendido.
    \item $\overline{p}_r$: potencia máxima del generador renovable $r \in \mathcal{R}$.
    \item $D_{n,t}$: demanda en el nodo $n \in N$ en la hora $t \in T$.
    \item $\overline{F}_l$: flujo máximo permitido por la línea de transmisión $l \in L$.
    \item $w^{RES}_{r,t}$: perfil de generación del generador renovable $r \in \mathcal{R}$ en la hora $t \in T$.
    \item $VoLL$: valor de la demanda no suministrada.
    \item $X_l$: reactancia de la línea $l \in L$.
\end{itemize}

\subsection*{Variables de decisión}

\begin{itemize}
    \item $P_{g,t}$: potencia generada por cada generador convencional $g \in G$ en la hora $t \in T$.
    \item $P^{RES}_{r,t}$: potencia generada por cada generador renovable $r \in \mathcal{R}$ utilizada para abastecer la demanda en la hora $t \in T$.
    \item $\Theta_{n,t}$: ángulo de fase en el nodo $n \in N$ durante la hora $t \in T$.
    \item $\Delta P_{n,t}$: potencia no suministrada en el nodo $n \in N$ durante el período $t \in T$.
    \item $F_{l,t}$: flujo de potencia por la línea $l \in L$ en la hora $t \in T$.
    \item $X^E_{g,t}$: variable binaria que indica el estado encendido/apagado del generador $g \in G$ en la hora $t \in T$.
    \item $X^{ON}_{g,t}$: variable binaria que indica si el generador convencional $g \in G$ se encendió en la hora $t \in T$.
    \item $X^{OFF}_{g,t}$: variable binaria que indica si el generador convencional $g \in G$ se apagó en la hora $t \in T$.
\end{itemize}

\subsection*{Función objetivo}

\begin{equation}
\begin{aligned}
\min \quad 
& \sum_{t \in T}\sum_{g \in G} C_g P_{g,t}
+ \sum_{t \in T}\sum_{n \in N} VoLL \, \Delta P_{n,t} 
 + \sum_{t \in T}\sum_{g \in G}
\left(
C^E_g X^{ON}_{g,t}
+ C^A_g X^{OFF}_{g,t}
\right)
\end{aligned}
\label{eq:objective}
\end{equation}

\subsection*{Restricciones}

\subsubsection*{Restricción de balance}

La potencia generada por los generadores convencionales y renovables en cada nodo $n \in N$, más los flujos entrantes y menos los flujos salientes, debe ser igual a la demanda suministrada en cada hora $t \in T$.

\begin{equation}
\sum_{g \in G_n} P_{g,t}
+ \sum_{r \in R_n} P^{RES}_{r,t}
+ \sum_{\substack{l \in L \\ to(l)=n}} F_{l,t}
- \sum_{\substack{l \in L \\ from(l)=n}} F_{l,t}
=
D_{n,t} - \Delta P_{n,t},
\quad \forall n \in N,\; t \in T.
\label{eq:power_balance}
\end{equation}

\subsubsection*{Restricción de mínimo y máximo técnico}

La potencia generada por cada generador convencional $g \in G$ no puede superar su potencia máxima ni ser inferior a su mínimo técnico cuando el generador está encendido.

\begin{equation}
X^E_{g,t} \underline{P}_g
\leq
P_{g,t}
\leq
\overline{P}_g X^E_{g,t},
\quad \forall g \in G,\; t \in T.
\label{eq:generation_limits}
\end{equation}

\subsubsection*{Restricción de encendido y apagado}

El estado binario de cada generador debe ser consistente con las variables binarias de encendido y apagado.

\begin{equation}
X^E_{g,t} - X^E_{g,t-1}
=
X^{ON}_{g,t} - X^{OFF}_{g,t},
\quad \forall g \in G,\; t \in \{2,\ldots,12\}.
\label{eq:commitment_transition}
\end{equation}

\subsubsection*{Restricción de borde de encendido y apagado}

El estado binario de cada generador en el inicio del período debe ser consistente con las variables de encendido y apagado iniciales.

\begin{equation}
X^E_{g,1}
=
X^{ON}_{g,1} - X^{OFF}_{g,1},
\quad \forall g \in G.
\label{eq:initial_commitment}
\end{equation}

\subsubsection*{Restricción de generación renovable}

La potencia generada por los generadores renovables no puede superar el perfil de generación disponible multiplicado por su capacidad máxima.

\begin{equation}
0
\leq
P^{RES}_{r,t}
\leq
w^{RES}_{r,t} \overline{p}_r,
\quad \forall r \in \mathcal{R},\; t \in T.
\label{eq:renewable_generation}
\end{equation}

\subsubsection*{Restricción de capacidad de línea}

El flujo de potencia en ambos sentidos para cada línea de transmisión $l \in L$ no puede superar el flujo máximo permitido.

\begin{equation}
-\overline{F}_l
\leq
F_{l,t}
\leq
\overline{F}_l,
\quad \forall l \in L,\; t \in T.
\label{eq:line_capacity}
\end{equation}

\subsubsection*{Restricción de energía no suministrada}

La potencia no suministrada en el nodo $n \in N$ durante el período $t \in T$ no puede superar la demanda nodal correspondiente.

\begin{equation}
0
\leq
\Delta P_{n,t}
\leq
D_{n,t},
\quad \forall n \in N,\; t \in T.
\label{eq:load_shedding}
\end{equation}

\subsubsection*{Restricción de flujo DC}

El flujo por cada línea se define como la diferencia angular entre los nodos conectados por la línea, dividida por su respectiva reactancia.

\begin{equation}
F_{l,t}
=
\frac{
\Theta_{from(l),t} - \Theta_{to(l),t}
}{
X_l
},
\quad \forall l \in L,\; t \in T.
\label{eq:dc_power_flow}
\end{equation}

\subsubsection*{Nodo de referencia}

Se fija un nodo de referencia del sistema eléctrico mediante la siguiente igualdad:

\begin{equation}
\Theta_{1,t} = 0,
\quad \forall t \in T.
\label{eq:reference_bus}
\end{equation}

\subsubsection*{Naturaleza de las variables}

\begin{equation}
P_{g,t} \geq 0,
\quad
P^{RES}_{r,t} \geq 0,
\quad
\Delta P_{n,t} \geq 0,
\quad \forall g \in G,\; r \in \mathcal{R},\; n \in N,\; t \in T.
\label{eq:continuous_variables}
\end{equation}

\begin{equation}
X^E_{g,t}, X^{ON}_{g,t}, X^{OFF}_{g,t}
\in \{0,1\},
\quad \forall g \in G,\; t \in T.
\label{eq:binary_variables}
\end{equation}
\end{comment}

\begin{comment}
Con el fin de garantizar la seguridad y confiabilidad del sistema eléctrico, se solicita modificar el modelo determinista anterior para incorporar unrequerimiento de \emph{reservas}. La idea es disponer de un margen de capacidad adicional que permita absorber desviaciones de la demanda o de la generación solar respecto a su pronóstico.

Un enfoque simple consiste en agregar una restricción que imponga que la
capacidad disponible del sistema en cada período supere a la demanda en una
fracción $\alpha \geq 0$, donde $\alpha$ representa el nivel de reserva exigido
(por ejemplo, $\alpha = 0.15$ para un $15\%$). Formule dicha restricción y
agréguela al modelo. Al hacerlo, considere lo siguiente:

\begin{enumerate}
    \item La restricción debe imponerse para cada período $t \in T$.
    \item Debe involucrar la capacidad de las unidades generadoras
    convencionales que se encuentran comprometidas, la generación solar
    disponible y la demanda del período.
    \item El parámetro $\alpha$ debe estar definido de modo que el caso
    $\alpha = 0$ sea \textbf{equivalente al escenario base sin reservas}.
    Verifique que su formulación cumple esta propiedad.
\end{enumerate}

Cabe destacar que este es un enfoque simplificado y demostrativo, cuyo
propósito es contrastarse con la formulación estocástica. En la práctica, un
modelo de reservas más riguroso define una variable de reserva explícita por
unidad generadora, acotada por la capacidad no despachada de cada unidad, y
distingue así la energía efectivamente despachada de la capacidad mantenida
como respaldo.
\end{comment}

%\subsubsection*{Enfoque A: Ponderación de la demanda en el balance}
%Este método no requiere añadir nuevas variables al problema. Consiste en modificar la ecuación de balance de potencia del sistema, indexada para cada período de tiempo, de modo que la demanda total a satisfacer sea penalizada o sobredimensionada por un factor $(1 + \alpha)$.

%\subsubsection*{Enfoque B: Variables de reserva por generador}
%Este método introduce flexibilidad al permitir que cada unidad decida cuánta reserva aportar. Se deben realizar las siguientes modificaciones:
%\begin{itemize}
%    \item Definir una nueva variable $R_{g,t}\geq0$ de decisión para cada generador (y por cada período de tiempo) que represente la potencia reservada para contingencias.
%    \item Incorporar una nueva restricción de seguridad que imponga que la suma de las reservas asignadas en todo el sistema sea mayor o igual a una fracción $\alpha$ de la demanda de ese período.
%    \item Modificar las restricciones de capacidad máxima de los generadores, asegurando que la combinación de la potencia despachada y la reserva asignada no supere la capacidad física de la unidad.
%\end{itemize}

\subsection*{B) Modelo estócastico}
En esta sección implementará el modelo estocástico de dos etapas. Las decisiones de compromiso de unidades, dadas por $u_{g,t}$, $v_{g,t}$ y $w_{g,t}$, corresponden a decisiones de primera etapa y, por lo tanto, no dependen del escenario. En cambio, las decisiones de despacho, tales como $p_{g,t,s}$, $p^{\text{sol}}_{t,s}$ y $\ell_{t,s}$, corresponden a decisiones de segunda etapa y pueden adaptarse a cada escenario. %El desglose de conjuntos, parámetros y variables para el modelo estocástico de dos etapas se muestra en la Tabla \ref{tab:notacion_est}.


%A medida que avance en la implementación, procure identificar claramente la diferencia entre decisiones tomadas antes y después de observar la incertidumbre.

\begin{comment}
\begin{table}[h]
\centering
\renewcommand{\arraystretch}{1.3}
\begin{tabular}{c l}
\toprule
\textbf{Símbolo} & \textbf{Descripción} \\
\midrule
\multicolumn{2}{l}{\textbf{\textit{Conjuntos e índices}}} \\
\midrule
$\mathcal{G}$ & Conjunto de unidades generadoras convencionales (índice $g$) \\
$\mathcal{T}$ & Conjunto de períodos de operación (índice $t$) \\
$\mathcal{S}$ & Conjunto de escenarios (índice $s$) \\
\midrule
\multicolumn{2}{l}{\textbf{\textit{Parámetros}}} \\
\midrule
$\pi_{s}$ & Probabilidad asociada al escenario $s$ \\
$D_{t,s}$ & Demanda eléctrica en el período $t$ bajo el escenario $s$ \\
$\overline{P}^{\text{sol}}_{t,s}$ & Generación solar máxima disponible en el período $t$ bajo el escenario $s$ \\
$c_{g}^{\text{var}}$ & Costo variable de generación de la unidad $g$ \\
$c_{g}^{\text{start}}$ & Costo de encendido (arranque) de la unidad $g$ \\
$\text{VOLL}$ & Costo de energía no suministrada (desprendimiento de carga) \\
\midrule
\multicolumn{2}{l}{\textbf{\textit{Variables de primera etapa}}} \\
\midrule
$u_{g,t}\in\{0,1\}$ & $=1$ si la unidad $g$ está encendida en el período $t$ \\
$v_{g,t}\in\{0,1\}$ & $=1$ si la unidad $g$ se enciende al inicio del período $t$ \\
$w_{g,t}\in\{0,1\}$ & $=1$ si la unidad $g$ se apaga al inicio del período $t$ \\
\midrule
\multicolumn{2}{l}{\textbf{\textit{Variables de segunda etapa}}} \\
\midrule
$p_{g,t,s}\geq 0$ & Potencia generada por la unidad $g$ en el período $t$ bajo el escenario $s$ \\
$p^{\text{sol}}_{t,s}\geq 0$ & Generación solar utilizada en el período $t$ bajo el escenario $s$ \\
$\ell_{t,s}\geq 0$ & Energía no suministrada en el período $t$ bajo el escenario $s$ \\
\bottomrule
\end{tabular}
\caption{Conjuntos, parámetros y variables del modelo estocástico de \textit{unit commitment} de dos etapas.}
\label{tab:notacion_est}
\end{table}
\end{comment}

%Al finalizar esta parte, identifique qué unidades son encendidas en cada período y cómo se reparte la generación entre las unidades convencionales y la generación solar disponible.

\subsection{Construcción de escenarios}

Posteriormente, incorpore incertidumbre en la demanda y en la generación solar mediante un conjunto finito de escenarios. Cada escenario representa una posible realización futura de las variables inciertas. Así, para cada escenario $s \in \mathcal{S}$, se consideran perfiles de demanda $D_{t,s}$ y generación solar disponible $\overline{P}^{\text{sol}}_{t,s}$. Estos escenarios pueden interpretarse como trayectorias alternativas del sistema durante el horizonte de planificación.

Revise cómo se construyen estos escenarios en el Notebook y discuta brevemente su impacto sobre las decisiones operacionales del sistema.

%\subsection{Formulación estocástica}

\subsection{Implementación computacional}

Complete o modifique el código entregado en el Jupyter Notebook para resolver el problema con reservas y luego estocástico. En particular, para el problema estocástico, verifique que:

\begin{enumerate}
\item Las variables binarias de compromiso no dependan del escenario.
\item Las variables de despacho sí dependan del escenario.
\item La función objetivo considere el costo esperado sobre todos los escenarios.
\item Las restricciones de balance de potencia se impongan para cada período y escenario.
\item La disponibilidad solar y la demanda sean correctamente indexadas por escenario.
\end{enumerate}

\subsection{Análisis de resultados}

Una vez resuelto el modelo, analice los resultados obtenidos. Algunas preguntas orientadoras son:

\begin{enumerate}
\item ¿Qué unidades generadoras son encendidas en cada período?
\item ¿Cómo cambia el despacho convencional entre escenarios?
\item ¿En qué períodos la incertidumbre de generación solar tiene mayor impacto?
\item ¿Existe energía no suministrada? En caso afirmativo, ¿en qué escenarios ocurre?
\item ¿Cómo se compara la solución estocástica con una solución determinista basada en valores esperados?
\item ¿Qué interpretación tiene el costo esperado obtenido por el modelo?
\end{enumerate}

\section{Discusión de resultados}

Responda al menos las siguientes preguntas:

\begin{enumerate}

\item Explique brevemente cuál es la diferencia entre una decisión de primera etapa y una decisión de segunda etapa en el contexto del problema implementado.

\item Identifique cuáles variables del modelo corresponden a decisiones de primera etapa y cuáles corresponden a decisiones de segunda etapa.

\item Explique por qué las variables binarias de compromiso de unidades no deberían depender del escenario.

\item Explique cómo se incorpora la incertidumbre en demanda y generación solar dentro del modelo.

\item Compare conceptualmente la formulación determinista con la formulación estocástica.

\item Analice los resultados obtenidos en el Notebook. Comente el patrón de encendido de unidades, el uso de generación solar y la presencia o ausencia de energía no suministrada.

\item Discuta qué ocurriría con el tamaño del problema si aumentara el número de escenarios considerados.

\item Compare sus resultados con el \textit{Jupyter Notebook} disponible en \href{https://github.com/paoptlab/Apunte-Optimizacion-A2IC}{link}.

\end{enumerate}


%\begin{comment}
\section{(Extra) Optimización a gran escala}

El equivalente determinista del problema estocástico puede crecer rápidamente a medida que aumenta el número de escenarios. Esto se debe a que las variables y restricciones de segunda etapa deben replicarse para cada escenario. En consecuencia, el tamaño del problema puede volverse computacionalmente desafiante.

Una estrategia clásica para abordar este tipo de problemas es la \textit{descomposición de Benders}, también conocida en programación estocástica como \textit{L-shaped method}. La idea central consiste en separar el problema en un problema maestro, que contiene las decisiones de primera etapa, y un conjunto de subproblemas, uno por cada escenario, asociados a las decisiones de segunda etapa. Para mayor detalle sobre la formulación,  consultar el complemento de la actividad publicado en material docente. 

Una vez estudiada la formulación, discuta cómo podría aplicarse esta descomposición al problema de UC estocástico. En particular, responda:

\begin{enumerate}
\item ¿Qué variables deberían quedar en el problema maestro?
\item ¿Qué variables deberían quedar en los subproblemas por escenario?
\item ¿Por qué la descomposición puede ser útil cuando aumenta el número de escenarios?
\item ¿Cómo deberían comportarse las cotas sobre el valor óptimo del problema ($LB$ y $UB$) a medida que se resuelven las iteraciones del algoritmo?
\end{enumerate}
%\end{comment}




% FIN DEL DOCUMENTO
\end{document}